% Propietario conservado: /Users/ruben/Documents/New project/output/EXTREMA_MEDIA_RAZON_RECIPROCIDAD_ES_EN_20260918/ARTICULO/ES/source/nuclear/02.tex
\chapter{Conservation des observables et registre complet}
\label{x_reciprocidad:nuclear:2}
\section{Bilan d’une observable conservée}

Soient \(\mathcal H,\mathcal K\) les réalisations d’un espace de lecture et de son extension enrichie. Soient \(J:\mathcal H\to\mathcal K\) une isométrie, \(P=JJ^*\), \(U:\mathcal K\to\mathcal K\) un transport borné et \(\widetilde A\) une observable bornée autoadjointe. Supposons

\[
U^*\widetilde A U=\widetilde A,\qquad [\widetilde A,P]=0.
\]

On définit

\[
A=J^*\widetilde A J,\qquad T=J^*UJ,
\qquad \eta=(I-P)UJ.
\]

\begin{proposition}
\label{x_reciprocidad:nuc02:balance}
La charge se décompose exactement entre lecture et complément :

\begin{equation}
\boxed{A=T^*AT+\eta^*\widetilde A\eta.}\label{x_reciprocidad:nuc02:eq:1}
\end{equation}
\end{proposition}

\begin{proof}
L’identité \(UJ=JT+\eta\) est une décomposition orthogonale. Lors du développement de \(J^*U^*\widetilde A UJ\), les termes croisés sont nuls parce que \(\widetilde A\) réduit les sous-espaces \(P\mathcal K\) et \((I-P)\mathcal K\). Le terme de gauche est \(A\) ; les deux termes diagonaux sont ceux de \eqref{x_reciprocidad:nuc02:eq:1}. Cela prouve également l’identité des formes sesquilinéaires pour deux états distincts. Si \(\widetilde A\geq0\), les deux termes sont positifs. Pour \(\widetilde A=I\), l’hypothèse de conservation est \(U^*U=I\).
\end{proof}

Si la lecture mélange des secteurs de l’observable, le bilan conserve en outre

\[
T^*J^*\widetilde A\eta+\eta^*\widetilde A JT.
\]

Ces termes représentent l’interférence entre les deux résidences de la charge. La condition de réduction écrite ci-dessus garantit leur annulation pour toute la décomposition. Dans des réalisations particulières, ils peuvent aussi s’annuler pour d’autres raisons, par exemple lorsque \(\eta=0\).

\section{Spécialisation à la connexion nonadique}

Soit \(C\) la réalisation unitaire du transport de mémoire sur \(\mathcal H\). Dans \(\mathcal H^9\), le transport des neuf phases et leur inclusion uniforme sont

\[
S_C(v_0,\ldots,v_8)=(Cv_8,v_0,\ldots,v_7),
\qquad Jv=\tfrac13(v,\ldots,v).
\]

Le facteur \(1/3\) résulte de la normalisation de neuf copies. La compression et son complément sont

\[
T=\frac{8I+C}{9},\qquad
\eta v=\frac1{27}\bigl(8(C-I)v,-(C-I)v,\ldots,-(C-I)v\bigr).
\]

Si \(A=A^*\) est borné et \([A,C]=0\), l’observable \(\widetilde A=I_9\otimes A\) satisfait les hypothèses de la proposition~\ref{x_reciprocidad:nuc02:balance}. Par conséquent,

\begin{equation}
\boxed{A=T^*AT+\frac8{81}(I-C)^*A(I-C).}\label{x_reciprocidad:nuc02:eq:2}
\end{equation}

Les coefficients proviennent de la répartition entre les neuf composantes : \(8^2+8=72\), divisé par \(27^2\), donne \(8/81\). L’identité est le bilan sesquilinéaire déjà documenté dans l’article spectral, exprimé ici pour une observable générale.

L’itération télescopique fournit, pour chaque \(N\geq1\),

\begin{equation}
A=T^{*N}AT^N+
\sum_{j=0}^{N-1}T^{*j}\eta^*\widetilde A\eta T^j.\label{x_reciprocidad:nuc02:eq:3}
\end{equation}

Pour \(A=I\), l’application

\[
V_Nv=(T^Nv,\eta v,\eta Tv,\ldots,\eta T^{N-1}v)
\]

est isométrique. L’étape \(N\to N+1\) décompose exclusivement le terminal \(T^Nv\) en sa nouvelle lecture et son complément ; elle maintient tous les registres précédents. Ainsi, la conservation se vérifie à toute profondeur.

Distinction des opérations. On a \(S_C^9=I_9\otimes C\). L’itération \(T^N\) décrit des compressions successives avec stockage de leurs compléments. Le transport enrichi sans compressions intermédiaires est \(S_C^N\). Les deux opérations ont des domaines définis et des fonctions différentes ; la preuve conserve cette distinction.

\section{Secteur terminal et registre à profondeur illimitée}

Soit \(P_{\rm fix}\) la projection orthogonale sur \(\ker(I-C)\).

\begin{theorem}
Pour tout transport unitaire \(C\), la compression nonadique vérifie

\begin{equation}
\underset{N\to\infty}{\operatorname{s-lim}}T^N=P_{\rm fix},\qquad
\boxed{I=P_{\rm fix}+\sum_{j=0}^{\infty}T^{*j}\eta^*\eta T^j.}\label{x_reciprocidad:nuc02:eq:4}
\end{equation}

La série converge dans la topologie forte. Pour chaque observable bornée autoadjointe \(A\) qui commute avec \(C\), on obtient en outre

\begin{equation}
\boxed{A=P_{\rm fix}AP_{\rm fix}
+\sum_{j=0}^{\infty}T^{*j}\eta^*(I_9\otimes A)\eta T^j.}\label{x_reciprocidad:nuc02:eq:5}
\end{equation}
\end{theorem}

\begin{proof}
Dans la représentation spectrale de l’opérateur unitaire \(C\), la compression correspond à \(t(z)=(8+z)/9\). Sur le cercle unité,

\[
|t(z)|^2=1-\frac8{81}|1-z|^2.
\]

Ainsi \(t(z)^N\to0\) pour \(z\ne1\), et \(t(1)^N=1\). La convergence dominée par rapport à chaque mesure spectrale vectorielle prouve \(T^N\to P_{\rm fix}\) fortement. Le même argument vaut pour \(T^{*N}\). Comme \(A\) est borné, \(T^{*N}AT^N\to P_{\rm fix}AP_{\rm fix}\) fortement. La limite de \eqref{x_reciprocidad:nuc02:eq:3} donne \eqref{x_reciprocidad:nuc02:eq:5}, et sa spécialisation \(A=I\) donne \eqref{x_reciprocidad:nuc02:eq:4}.
\end{proof}

En particulier,

\[
V_\infty v=(P_{\rm fix}v,\eta v,\eta Tv,\eta T^2v,\ldots)
\]

est une isométrie. Elle conserve le produit scalaire complet, en plus de la norme.

\begin{corollary}[Mémoire bilatérale]
Pour le décalage \(Ce_n=e_{n+1}\) dans \(\ell^2(\mathbb Z)\), \(P_{\rm fix}=0\) : une suite fixe serait constante et sa seule réalisation de carré sommable est zéro. En conséquence, toute la norme initiale est distribuée dans le registre des compléments. Cela précise le comportement illimité de la réalisation bilatérale de mémoire présente dans le corpus.
\end{corollary}

\begin{corollary}[Calendrier fini]
Dans un cycle de \(d\) positions, la suite uniforme détermine un sous-espace fixe unidimensionnel. Ce secteur reste terminal et le complément accumule le reste. Le calendrier fini et le registre bilatéral possèdent donc des terminaux différents.
\end{corollary}

Dans le cas bilatéral, la convergence forte peut coexister avec \(\|T^N\|=1\) pour tout \(N\). La démonstration ne suppose pas de taux uniforme de contraction. Elle distingue également cette compression du mode \((5/9)C\), dont la contraction géométrique est uniforme et étudiée séparément dans les sources.

\section{Conservation de la symétrie et transport de carte
}

Soit \(R\) unitaire et définissons \(C'=RCR^*\), \(A'=RAR^*\). La formule explicite de la compression donne

\begin{equation}
T(C')R=RT(C),\qquad
\eta(C')R=R^{\oplus9}\eta(C).\label{x_reciprocidad:nuc02:eq:6}
\end{equation}

\begin{proof}
La première égalité s’obtient en distribuant \((8I+RCR^*)R/9\). La seconde résulte de \((C'-I)R=R(C-I)\) dans chacune des neuf composantes de \(\eta\).
\end{proof}

L’itération de \eqref{x_reciprocidad:nuc02:eq:6} entrelace les registres finis. De plus, \(P_{\rm fix}(C')=RP_{\rm fix}(C)R^*\), par transport du noyau de \(I-C\), de sorte que le registre illimité est également entrelacé. Ainsi, le bilan des observables est covariant lorsque la carte, l’opérateur et l’observable sont transportés conjointement. Les symétries qui commutent avec \(C\) agissent en outre dans une même carte fixe. Cette distinction permet d’incorporer les transformations exceptionnelles avec leurs repères, leurs drapeaux et leurs orientations effectivement transportés.

La composition de \cref{x_reciprocidad:lem:k-hadamard,x_reciprocidad:exc:entrelazador,x_reciprocidad:exc:isometria} utilise la matrice de Hadamard du registre, \(H_{12}^2=4I\), et l’incidence \(\mathcal I\) des douze positions dans les 132 hexades :

\[
J_H=H_{12}/2,\qquad
\mathcal I^*\mathcal I=36I+30\mathbf1\mathbf1^*,\qquad
U_W=\mathcal I/6\big|_{\mathbf1^\perp}.
\]

Les deux opérateurs normalisés sont isométriques dans les domaines indiqués. La composition \(B=U_WJ_H\) sur \(W=J_H(\mathbf1^\perp)\) conserve les produits scalaires et entrelace la représentation transportée de \(M_{12}\) avec l’action sur les hexades. Le registre entier \(K\) conserve en outre l’inversion entière \(H_{12}/4\) sur son réseau de congruences ; cette inversion et la normalisation isométrique ont des fonctions différentes. La démonstration, le prolongement APP–Witt–Leech et leurs domaines sont développés ici.

Dans cette composition, la conservation de l’incidence a un contenu opératoriel concret. Son extension par d’autres transports utilise \eqref{x_reciprocidad:nuc02:eq:6} et les entrelaceurs effectifs des sources, en respectant leurs groupes et domaines respectifs.

\subsection{Transport intégral des douze coordonnées de K}

La séparation du sous-espace centré peut être complétée en conservant aussi sa moyenne. Soient \(P_0=I-\mathbf1\mathbf1^*/12\), \(\mu(k)=\mathbf1^*k/12\) et

\[
F(k)=\left(\mu(k),\frac16\mathcal I P_0k\right)\in\mathbb R\oplus E_{\rm inc}.
\]

Nous munissons le codomaine de \(\|(\mu,y)\|^2=12|\mu|^2+\|y\|^2\). Le Gram d’incidence démontre

\begin{equation}
\|F(k)\|^2=\|k\|^2,\qquad
\boxed{F^{-1}(\mu,y)=\mu\mathbf1+\frac16P_0\mathcal I^*y.}\label{x_reciprocidad:nuc02:eq:6a}
\end{equation}

En effet, \(\mathcal I^*\mathcal I P_0=36P_0\), de sorte que la formule inverse récupère \(P_0k\), et le premier terme récupère sa moyenne. L’action de \(M_{12}\) est triviale sur la moyenne et permute les hexades ; c’est pourquoi \(F\) est équivariant. Composer \(FJ_H\) transporte les coordonnées signées normalisées. Ainsi, les douze coordonnées du registre sont conservées, et pas seulement sa partie centrée. Le domaine de cette reconstruction est le registre dodécaphasique ; la mémoire complète de l’histoire maintient ses autres coordonnées et son prolongement.

\subsection{Spécialisation à l’opérateur exceptionnel d’ordre trois}

La construction exceptionnelle produit l’opérateur orthogonal \(g_c\) sur la réalisation du voisin de Leech, au moyen des douze fibres APP–Witt de type \(A_2\), et démontre \(g_c^3=I\) et \(\ker(I-g_c)=0\). Sa construction, son entrelacement et sa préservation réticulaire sont localisés dans \cref{x_reciprocidad:nuclear:fibras-app-witt,x_reciprocidad:nuc04:c3}.

Nous construisons maintenant, comme composition opératoire de cette sortie avec la carte à neuf phases,

\[
T_3=\frac{8I+g_c}{9},\qquad
\eta_3=(I-JJ^*)S_{g_c}J.
\]

\begin{corollary}
Cette composition vérifie

\begin{equation}
\boxed{T_3^*T_3=\frac{19}{27}I,\qquad
\eta_3^*\eta_3=\frac8{27}I,\qquad
\|T_3^Nv\|^2=\left(\frac{19}{27}\right)^N\|v\|^2.}\label{x_reciprocidad:nuc02:eq:6b}
\end{equation}
\end{corollary}

\begin{proof}
Le vecteur \((I+g_c+g_c^2)v\) est fixe par \(g_c\) ; l’absence de vecteurs fixes implique \(I+g_c+g_c^2=0\). Comme \(g_c^*=g_c^2\), nous avons \(g_c+g_c^*=-I\). Par substitution,

\[
T_3^*T_3=\frac{65I+8(g_c+g_c^*)}{81}
=\frac{57}{81}I=\frac{19}{27}I.
\]

Le bilan \eqref{x_reciprocidad:nuc02:eq:2} avec \(A=I\) donne le second coefficient. L’application réitérée de l’égalité des normes donne la puissance pour tout \(N\).
\end{proof}

Le registre illimité est isométrique et son terminal tend vers zéro au taux explicite de \eqref{x_reciprocidad:nuc02:eq:6b}. Ici, la structure exceptionnelle détermine les deux coefficients complémentaires ; tous deux proviennent de la relation polynomiale de \(g_c\) et des neuf phases. C’est une composition mathématique développée dans cette recherche. Sa définition n’identifie pas \(g_c\) à \(\Gamma_9\) et n’affirme pas que \(S_{g_c}\) soit l’opérateur temporel canonique du TPK : elle conserve la distinction entre prolongement exceptionnel et holonomie temporelle tout en les composant explicitement.

\section{Raffinement des histoires et compatibilité des bilans}

Pour les histoires enrichies de profondeur \(n\), soit \(p_h(\varepsilon)\) la probabilité d’une extension survivante, normalisée par \(\sum_\varepsilon p_h(\varepsilon)=1\). Le raffinement canonique sur les bases d’histoires est

\[
R_n\delta_h=\sum_{\varepsilon\in\operatorname{Out}(h)}
\sqrt{p_h(\varepsilon)}\,\delta_{h\varepsilon}.
\]

Les extensions d’histoires différentes ont des préfixes différents et des supports disjoints. C’est pourquoi \(R_n^*R_n=I\). C’est la réalisation isométrique de la cohérence projective de la mesure : les coefficients proviennent des poids de survie déjà construits.

Pour transporter \eqref{x_reciprocidad:nuc02:eq:1} entre profondeurs, on conserve conjointement les inclusions, les projections, les transports et les observables. Avec les isométries de raffinement \(R_n,\widetilde R_n\), les carrés nécessaires sont

\[
J_{n+1}R_n=\widetilde R_nJ_n,\quad
P_{n+1}\widetilde R_n=\widetilde R_nP_n,\quad
U_{n+1}\widetilde R_n=\widetilde R_nU_n,\quad
\widetilde A_{n+1}\widetilde R_n=\widetilde R_n\widetilde A_n.
\]

Ces carrés impliquent \(T_{n+1}R_n=R_nT_n\) et \(\eta_{n+1}R_n=\widetilde R_n\eta_n\). Avec des bornes uniformes, les identités passent à la limite inductive des réalisations. Les histoires, en revanche, forment le système projectif des préfixes ; les directions opposées des deux types d’applications sont distinguées. Les preuves de mesure, de matrice de Gram et de passage à la limite sont incluses.

\section{Spécialisation variationnelle : action et orientation}

L’action elliptique de \cref{x_reciprocidad:iv:ellipse:section,x_reciprocidad:nuclear:7} fournit une spécialisation physique du principe, avec les paramètres construits dans ses étapes HMT. Dans ses coordonnées canoniques,

\[
H_\eta(Q,P)=\frac{\omega}{2}
\left(e^{-\eta}Q^2+e^\eta P^2\right),\qquad
I_\eta=H_\eta/\omega.
\]

La transformation canonique \(Q=e^{\eta/2}q\), \(P=e^{-\eta/2}p\) transforme \(I_\eta\) en \((q^2+p^2)/2\). Les rotations de \((q,p)\), transportées vers \((Q,P)\), constituent une symétrie continue. Pour le lagrangien \(L=P\dot Q-\omega I_\eta\), sa charge de Noether est \(I_\eta\). Sur le contour de la source \(I_\eta=\hbar\), l’action orientée est \(\oint P\,dQ=2\pi\hbar=h\), avec l’orientation de l’orbite hamiltonienne. La preuve distingue la transformation symplectique de la transformation antisymplectique : elles peuvent conserver la même énergie et agir avec des signes différents sur l’action orientée. La démonstration variationnelle complète est présentée dans \cref{x_reciprocidad:nuclear:7}.

\subsection{Mise à jour de la forme et terme de connexion}

La matrice de forme \(M_\eta=\operatorname{diag}(e^{\eta/2},e^{-\eta/2})\) permet de développer un transport entre ellipses de paramètres différents. Nous fixons \(J_2=\begin{pmatrix}0&-1\\1&0\end{pmatrix}\) et \(R(\theta)=\exp(-\theta J_2)\), la rotation hamiltonienne dans le plan \((q,p)\). Le transport

\[
z'=M_{\eta'}R(\delta\theta)M_\eta^{-1}z,
\qquad z=(Q,P)^\mathsf T,
\]

est symplectique et satisfait exactement \(I_{\eta'}(z')=I_\eta(z)\). La preuve consiste à envoyer les deux états dans leurs coordonnées normalisées et à utiliser la conservation de \(q^2+p^2\) par la rotation. Il conserve l’action tout en changeant la forme.

Pour une réalisation différentiable \(\eta(t)\), d’avancée angulaire \(\dot\theta=\omega\), la différentiation de ce transport produit

\[
\dot Q=\omega e^\eta P+\tfrac12\dot\eta Q,\qquad
\dot P=-\omega e^{-\eta}Q-\tfrac12\dot\eta P.
\]

Son générateur hamiltonien, obtenu à partir de ces deux équations, est

\begin{equation}
\boxed{H_{\rm trans}=\omega I_\eta+\frac{\dot\eta}{2}QP.}\label{x_reciprocidad:nuc02:eq:7}
\end{equation}

Le second terme est le terme de connexion de la forme variable. Avec la convention \(\{f,g\}=\partial_Qf\partial_Pg-\partial_Pf\partial_Qg\),

\[
\partial_tI_\eta=\frac{\dot\eta}{2}(-e^{-\eta}Q^2+e^\eta P^2),
\quad
\left\{I_\eta,\frac{\dot\eta}{2}QP\right\}
=\frac{\dot\eta}{2}(e^{-\eta}Q^2-e^\eta P^2).
\]

Les deux contributions s’annulent et \(dI_\eta/dt=0\). Pour le transport prescrit, \eqref{x_reciprocidad:nuc02:eq:7} est unique à une fonction du temps près. La réalisation différentiable est une spécialisation déclarée du transport entre formes ; le résultat discret précédent vaut directement entre deux paramètres engendrés. La mise à jour TPK spécifique conserve sa propre règle de sélection de ces paramètres et de leurs temps.

On obtient ainsi une connexion précise : un observable conservé par la dynamique complète se distribue entre lecture et mémoire au moyen de \eqref{x_reciprocidad:nuc02:eq:1}. Lorsque cet observable est le générateur d’une symétrie continue de l’action construite, la quantité distribuée est une charge de Noether. La norme, la charge d’une symétrie et la fonctionnelle énergétique sont reliées par leurs opérateurs spécifiques.

Une spécialisation bornée effective s’obtient à partir du calendrier de 108 états. Son hamiltonien réalisé est \(H_{\rm clk}=\hbar\omega_0N_{108}\), avec \(N_{108}=\sum_{k=0}^{107}kP_k\) et \(P_k\) ses projecteurs de Fourier. L’action de l’horloge a pour charge \(\hbar\langle\psi,N_{108}\psi\rangle=E_{\rm clk}/\omega_0\). Les transports et les inclusions qui satisfont les hypothèses de \eqref{x_reciprocidad:nuc02:eq:1} distribuent cette charge entre lecture et complément. Ici, tous les opérateurs sont finis et bornés. L’action elliptique de \eqref{x_reciprocidad:nuc02:eq:7} a été prouvée dans sa réalisation classique symplectique ; une quantification avec des observables non bornées conserve en outre ses propres conditions de domaine.
